%% kapa: exam=1 round=37 year=2026 period=1 form=A subject=economics subject_ko=경제학원론 seq=35 number=75
%% source: Q-Net 감정평가사 기출문제 (cst00309) artlSeq=5258924 file=raw/1008/5258924/2026년 제37회 감정평가사 1차 1교시.pdf
%% conversion: kordoc 4.24.51 -> tools/gen1.py (auto); review=proofread-sonnet; flags=none
\begin{question}{75}
( )에 들어갈 내용으로 옳은 것은? (단, 전쟁 이후 인구증가율, 저축률, 감가상각률은 변하지 않았다고 가정한다.)
\begin{framedtext}
균제상태(steady state)에 있던 A국에서 전쟁으로 많은 노동인구가 감소했으나 공장 등 자본스톡의 피해는 없었다. 솔로우(Solow) 성장모형에 의하면 1인당 산출량은 전쟁 직후에 ( ㄱ ), 다시 균제상태에 이르면 전쟁 전과 비교하여 ( ㄴ ).
\end{framedtext}
\begin{choices}
  \item ㄱ: 증가하고, ㄴ: 같아진다
  \item ㄱ: 증가하고, ㄴ: 감소한다
  \item ㄱ: 감소하고, ㄴ: 같아진다
  \item ㄱ: 감소하고, ㄴ: 감소한다
  \item ㄱ: 변하지 않고, ㄴ: 같아진다
\end{choices}
\end{question}
\answer{1}
